\ifdefined\gkver\else\def\gkver{student}\fi
\documentclass[\gkver]{gaokaozhenti}
\begin{document}

\chapter{2004年全国理综Ⅰ}
%% paperType: 理综物理部分

\begin{enumerate}
\item
%% number: 14
%% paperName: 2004年全国理综Ⅰ第 14 题 6 分
%% typeId: 1
%% type: 单选题
%% chapter: 原子和原子核
%% point: 天然放射性
%% method: 无
%% score: 6
%% degree: 650
%% duplicateId: 0
%% body:
本题中用大写字母代表原子核。E 经 $\alpha$ 衰变成为 F，再经 $\beta$ 衰变成为 G，再经 $\alpha$ 衰变成为 H。上述系列衰变可记为下式：E$\xrightarrow{\alpha}$F$\xrightarrow{\beta}$G$\xrightarrow{\alpha}$H

另一系列衰变如下：P$\xrightarrow{\beta}$Q$\xrightarrow{\beta}$R$\xrightarrow{\alpha}$S

已知 P 是 F 的同位素，则\xzanswer{B}
\fourchoices[answer=B]
{Q 是 G 的同位素，R 是 H 的同位素}
{R 是 E 的同位素，S 是 F 的同位素}
{R 是 G 的同位素，S 是 H 的同位素}
{Q 是 E 的同位素，R 是 F 的同位素}

%% answer: B
%% memo:
\memoanswer{
	写出衰变通式 $\ce{^{A}_{Z}E -> ^{A-4}_{Z-2}F -> ^{A-4}_{Z-1}G -> ^{A-8}_{Z-3}H}$，因 P 是 F 的同位素，故有 $\ce{^{B}_{Z-2}P -> ^{B}_{Z-1}Q -> ^{B}_{Z}R -> ^{B}_{Z-2}S}$，比较可知 R 是 E 的同位素，S 是 F 的同位素，Q 是 G 的同位素，故 B 正确。
}

\item
%% number: 15
%% paperName: 2004年全国理综Ⅰ第 15 题 6 分
%% typeId: 1
%% type: 单选题
%% chapter: 运动和力
%% point: 牛顿定律的应用
%% method: 无
%% score: 6
%% degree: 550
%% duplicateId: 0
%% body:
如图所示，ad、bd、cd 是竖直面内三根固定的光滑细杆，a、b、c、d 位于同一圆周上，a 点为圆周的最高点，d 点为最低点。每根杆上都套着一个小滑环（图中未画出），三个滑环分别从 a、b、c 处释放（初速为 0），用 $t_{1}$、$t_{2}$、$t_{3}$ 依次表示各滑环到达 d 所用的时间，则\xzanswer{D}
\onepicture[h=4.4cm]{figs/15.png}
\fourchoices[answer=D]
{$t_{1}<t_{2}<t_{3}$}
{$t_{1}>t_{2}>t_{3}$}
{$t_{3}>t_{1}>t_{2}$}
{$t_{1}=t_{2}=t_{3}$}

%% answer: D
%% memo:
\memoanswer{
	设杆与水平面夹角为 $\theta$，圆的直径为 $d$，由牛顿第二定律得 $mg\sin\theta=ma$，解得 $a=g\sin\theta$，由运动学规律得 $s=\frac{1}{2}at^{2}$，又 $s=d\sin\theta$，解得 $t=\sqrt{\frac{2d}{g}}$，与 $\theta$ 无关，则 $t_{1}=t_{2}=t_{3}$，故 D 正确。
}

\item
%% number: 16
%% paperName: 2004年全国理综Ⅰ第 16 题 6 分
%% typeId: 1
%% type: 单选题
%% chapter: 气体性质
%% point: 阿伏伽德罗常数
%% method: 无
%% score: 6
%% degree: 550
%% duplicateId: 0
%% body:
若以 $\mu$ 表示水的摩尔质量，$\upsilon$ 表示在标准状态下水蒸气的摩尔体积，$\rho$ 为在标准状态下水蒸气的密度，$N_{A}$ 为阿佛加德罗常数，$m$、$\Delta$ 分别表示每个水分子的质量和体积，下面是四个关系式：

①$N_{A}=\frac{\upsilon\rho}{m}$\quad ②$\rho=\frac{\mu}{N_{A}\Delta}$\quad ③$m=\frac{\mu}{N_{A}}$\quad ④$\Delta=\frac{\upsilon}{N_{A}}$

其中\xzanswer{B}
\fourchoices[answer=B]
{①和②都是正确的}
{①和③都是正确的}
{③和④都是正确的}
{①和④都是正确的}

%% answer: B
%% memo:
\memoanswer{
	由微观量与宏观量的运算关系知：$N_{A}=\frac{\rho\upsilon}{m}$ 正确；关系式 $\rho=\frac{\mu}{N_{A}\Delta}$ 中的 $\frac{\mu}{N_{A}\Delta}$ 为水的密度，不等于水蒸气的密度，不正确；$m=\frac{\mu}{N_{A}}$ 正确；关系式 $\Delta=\frac{\upsilon}{N_{A}}$ 中的 $\frac{\upsilon}{N_{A}}$ 为每个水蒸气分子占据的空间，不正确，故 B 正确。
}

\item
%% number: 17
%% paperName: 2004年全国理综Ⅰ第 17 题 6 分
%% typeId: 1
%% type: 单选题
%% chapter: 机械振动
%% point: 振动图象
%% method: 无
%% score: 6
%% degree: 450
%% duplicateId: 0
%% body:
一列简谐横波沿 $x$ 轴负方向传播，图\ref{2004全国理综Ⅰ17a}是 $t=1\Us$ 时的波形图，图\ref{2004全国理综Ⅰ17b}是波中某振动质元位移随时间变化的振动图线（两图用同一时间起点），则图\ref{2004全国理综Ⅰ17b}可能是图\ref{2004全国理综Ⅰ17a}中哪个质元的振动图线？\xzanswer{A}
\twopicture[h=3.4cm, labelA=2004全国理综Ⅰ17a, labelB=2004全国理综Ⅰ17b]
{figs/17a.png}
{figs/17b.png}
\fourchoices[answer=A]
{$x=0$ 处的质元}
{$x=1\Um$ 处的质元}
{$x=2\Um$ 处的质元}
{$x=3\Um$ 处的质元}

%% answer: A
%% memo:
\memoanswer{
	当 $t=1\Us$ 时，由振动图线知，质点在平衡位置且向下运动，在波的图象中，由波向左传播可知 $x=0$ 处的质点位于平衡位置且向下运动，故 A 正确。
}

\item
%% number: 18
%% paperName: 2004年全国理综Ⅰ第 18 题 6 分
%% typeId: 1
%% type: 单选题
%% chapter: 电路
%% point: 串并联电路
%% method: 无
%% score: 6
%% degree: 650
%% duplicateId: 0
%% body:
图中电阻 $R_{1}$、$R_{2}$、$R_{3}$ 的阻值相等，电池的内阻不计。开关 K 接通后流过 $R_{2}$ 的电流是 K 接通前的\xzanswer{B}
\onepicture[h=3.4cm]{figs/18.png}
\fourchoices[answer=B]
{$\frac{1}{2}$}
{$\frac{2}{3}$}
{$\frac{1}{3}$}
{$\frac{1}{4}$}

%% answer: B
%% memo:
\memoanswer{
	开关 K 接通前，$R_{1}$ 和 $R_{2}$ 串联，$I_{1}=\frac{E}{2R}$；开关 K 接通后，$R_{2}$ 和 $R_{3}$ 并联，再与 $R_{1}$ 串联，$I_{2}'=\frac{1}{2}I_{2}=\frac{1}{2}\cdot\frac{E}{R+\frac{R}{2}}=\frac{E}{3R}$，故 $\frac{I_{2}'}{I_{1}}=\frac{2}{3}$，故 B 正确。
}

\item
%% number: 19
%% paperName: 2004年全国理综Ⅰ第 19 题 6 分
%% typeId: 1
%% type: 单选题
%% chapter: 光学
%% point: 光电效应
%% method: 无
%% score: 6
%% degree: 650
%% duplicateId: 0
%% body:
下表给出了一些金属材料的逸出功。

\begin{center}
\begin{tabular}{|c|c|c|c|c|c|}
\hline
材料 & 铯 & 钙 & 镁 & 铍 & 钛 \\
\hline
逸出功（$10^{-19}\UJ$） & 3.0 & 4.3 & 5.9 & 6.2 & 6.6 \\
\hline
\end{tabular}
\end{center}

现用波长为 $400\Unm$ 的单色光照射上述材料，能产生光电效应的材料最多有几种？（普朗克常量 $h=6.6\times10^{-34}\UJcdots$，光速 $c=3.0\times10^{8}\Ums$）\xzanswer{A}
\fourchoices[answer=A]
{2 种}
{3 种}
{4 种}
{5 种}

%% answer: A
%% memo:
\memoanswer{
	波长 $400\Unm$ 光子的能量 $E=h\frac{c}{\lambda}=6.6\times10^{-34}\times\frac{3.0\times10^{8}}{400\times10^{-9}}\UJ=4.95\times10^{-19}\UJ$，大于铯和钙的逸出功，故 A 正确。
}

\item
%% number: 20
%% paperName: 2004年全国理综Ⅰ第 20 题 6 分
%% typeId: 1
%% type: 单选题
%% chapter: 运动和力
%% point: 超重失重
%% method: 无
%% score: 6
%% degree: 650
%% duplicateId: 0
%% body:
下列哪个说法是正确的？\xzanswer{B}
\fourchoices[answer=B]
{游泳运动员仰卧在水面静止不动时处于失重状态}
{蹦床运动员在空中上升和下落过程中都处于失重状态}
{举重运动员在举起杠铃后不动的那段时间内处于超重状态}
{体操运动员双手握住单杠吊在空中不动时处于失重状态}

%% answer: B
%% memo:
\memoanswer{
	物体具有向下的加速度时，处于失重状态；具有向上的加速度时，处于超重状态。ACD 项的运动员均处于平衡状态，B 中运动员为失重状态，故 B 正确。
}

\item
%% number: 21
%% paperName: 2004年全国理综Ⅰ第 21 题 6 分
%% typeId: 1
%% type: 单选题
%% chapter: 光学
%% point: 光的折射
%% method: 无
%% score: 6
%% degree: 450
%% duplicateId: 0
%% body:
发出白光的细线光源 ab，长度为 $l_{0}$，竖直放置，上端 a 恰好在水面以下，如图。现考虑线光源 ab 发出的靠近水面法线（图中的虚线）的细光束经水面折射后所成的像，由于水对光有色散作用，若以 $l_{1}$ 表示红光成的像的长度，$l_{2}$ 表示蓝光成的像的长度，则\xzanswer{D}
\onepicture[h=3.8cm]{figs/21.png}
\fourchoices[answer=D]
{$l_{1}<l_{2}<l_{0}$}
{$l_{1}>l_{2}>l_{0}$}
{$l_{2}>l_{1}>l_{0}$}
{$l_{2}<l_{1}<l_{0}$}

%% answer: D
%% memo:
\memoanswer{
	如图所示，由于蓝光折射率比红光大，故同一点发出的光经水面折射后，蓝光比红光偏折角大，经反向延长后所成的虚像长度小，即 $l_{2}<l_{1}<l_{0}$，故 D 正确。

	\onepicture[w=3.6cm]{figs/21_an.png}
}

\item
%% number: 22
%% paperName: 2004年全国理综Ⅰ第 22 题 18 分
%% typeId: 4
%% type: 实验题
%% chapter: 电路
%% point: 测电动势与内阻
%% method: 无
%% score: 18
%% degree: 450
%% duplicateId: 0
%% body:
\begin{enumerate}
	\item
	图中给出的是用螺旋测微器测量一金属薄板厚度时的示数，此读数应为\_\_\_\_\_\_\_\_\_\_mm。

	\onepicture[align=right, w=5.0cm]{figs/22a.png}
	\item
	实验室内有一电压表 mV，量程为 $150\UmV$、内阻约为 $150\UO$。现要将其改装成量程为 $10\UmA$ 的电流表，并进行校准。为此，实验室提供如下器材：干电池 $E$（电动势为 $1.5\UV$），电阻箱 $R$，滑线变阻器 $R'$，电流表 A（有 $1.5\UmA$、$15\UmA$ 与 $150\UmA$ 三个量程）及开关 K。
	\onepicture[align=right, w=5.0cm]{figs/22b.png}

	（a）对电流表改装时必须知道电压表的内阻。可用图示的电路测量电压表 mV 的内阻。在既不损坏仪器又能使精确度尽可能高的条件下，电路中的电流表 A 应选用的量程是\_\_\_\_\_\_\_\_\_\_。若合上 K，调节滑线变阻器后测得电压表的读数为 $150\UmV$，电流表 A 的读数为 $1.05\UmA$，则电压表的内阻 $R_{mV}$ 为\_\_\_\_\_\_\_\_\_\_。（取三位有效数字）

	（b）在对改装成的电流表进行校准时，把 A 作为标准电流表，画出对改装成的电流表进行校准的电路原理图（滑线变阻器作限流使用），图中各元件要用题中给出的符号或字母标注。图中电阻箱的取值是\_\_\_\_\_\_\_\_（取三位有效数字），电流表 A 应选的量程是\_\_\_\_\_\_\_\_\_\_。
\end{enumerate}

%% answer: 
\jdanswer{
\begin{enumerate}
	\item
	$6.124\Umm$
	\item
	（a）$1.5\UmA$、$143\UO$；（b）如图所示、$16.8\UO$、$15\UmA$。

	\onepicture[w=5.0cm]{figs/22_an.png}
\end{enumerate}
}
%% memo:
\memoanswer{
	（1）由图知螺旋测微器的读数为 $6\Umm+12.4\times0.01\Umm=6.124\Umm$；

	（2）（a）电压表满偏电流 $I=\frac{U}{R}=\frac{150\UmV}{150\UO}=1.0\UmA$，所以选择 $1.5\UmA$ 量程的电流表，电压表的内阻为 $R_{mV}=\frac{150\UmV}{1.05\UmA}\approx142.9\UO$，题干要求保留三位有效数字，故 $R_{mV}\approx143\UO$；

	（b）改装成 $10\UmA$ 的电流表后，校准时应选用 $15\UmA$ 量程的电流表，并联电阻箱的阻值为 $R=\frac{150\times10^{-3}}{(10-1.05)\times10^{-3}}\UO=16.8\UO$。
}

\item
%% number: 23
%% paperName: 2004年全国理综Ⅰ第 23 题 16 分
%% typeId: 6
%% type: 计算题
%% chapter: 曲线运动
%% point: 斜抛运动
%% method: 无
%% score: 16
%% degree: 450
%% duplicateId: 0
%% body:
在勇气号火星探测器着陆的最后阶段，着陆器降落到火星表面上，再经过多次弹跳才停下来。假设着陆器第一次落到火星表面弹起后，到达最高点时高度为 $h$，速度方向是水平的，速度大小为 $v_{0}$，求它第二次落到火星表面时速度的大小，计算时不计火星大气阻力。已知火星的一个卫星的圆轨道的半径为 $r$，周期为 $T$，火星可视为半径为 $r_{0}$ 的均匀球体。
%% answer: 
\jdanswer{
\[ v=\sqrt{\frac{8\pi^{2}hr^{3}}{T^{2}r_{0}^{2}}+v_{0}^{2}} \]
}
%% memo:
\memoanswer{
	以 $g'$ 表示火星表面附近的重力加速度，$M$ 表示火星的质量，$m$ 表示火星的卫星的质量，$m'$ 表示火星表面处某一物体的质量，由万有引力定律和牛顿第二定律得

	\[ G\frac{Mm'}{r_{0}^{2}}=m'g',\quad G\frac{Mm}{r^{2}}=m\left(\frac{2\pi}{T}\right)^{2}r, \]

	设着陆器第二次落到火星表面时的速度为 $v$，它的竖直分量为 $v_{1}$，水平分量仍为 $v_{0}$，有

	\[ v_{1}^{2}=2g'h,\quad v=\sqrt{v_{1}^{2}+v_{0}^{2}}, \]

	联立解得 $v=\sqrt{\frac{8\pi^{2}hr^{3}}{T^{2}r_{0}^{2}}+v_{0}^{2}}$。
}

\item
%% number: 24
%% paperName: 2004年全国理综Ⅰ第 24 题 18 分
%% typeId: 6
%% type: 计算题
%% chapter: 电磁感应
%% point: 电磁感应能量分析
%% method: 无
%% score: 18
%% degree: 450
%% duplicateId: 0
%% body:
图中 $a_{1}b_{1}c_{1}d_{1}$ 和 $a_{2}b_{2}c_{2}d_{2}$ 为在同一竖直平面内的金属导轨，处在磁感强度 $B$ 的匀强磁场中，磁场方向垂直导轨所在平面（纸面）向里。导轨的 $a_{1}b_{1}$ 段与 $a_{2}b_{2}$ 段是竖直的，距离为 $l_{1}$；$c_{1}d_{1}$ 段与 $c_{2}d_{2}$ 段也是竖直的，距离为 $l_{2}$。$x_{1}y_{1}$ 与 $x_{2}y_{2}$ 为两根用不可伸长的绝缘轻线相连的金属细杆，质量分别为 $m_{1}$、$m_{2}$，它们都垂直于导轨并与导轨保持光滑接触。两杆与导轨构成的回路的总电阻为 $R$。$F$ 为作用于金属杆 $x_{1}y_{1}$ 上竖直向上的恒力。已知两杆运动到图示位置时，已匀速向上运动，求此时作用于两杆的重力的功率的大小和回路电阻上的热功率。
\onepicture[align=right, h=5.5cm]{figs/24.png}
%% answer: 
\jdanswer{
\[ P=\frac{F-(m_{1}+m_{2})g}{B^{2}(l_{2}-l_{1})^{2}}R(m_{1}+m_{2})g,\quad Q=\left[\frac{F-(m_{1}+m_{2})g}{B(l_{2}-l_{1})}\right]^{2}R \]
}
%% memo:
\memoanswer{
	设杆向上运动的速度为 $v$，因杆的运动，两杆与导轨构成的回路的面积减少，从而磁通量也减少。由法拉第电磁感应定律，回路中的感应电动势的大小为

	\[ E=B(l_{2}-l_{1})v \tag{①} \]

	回路中的感应电流为

	\[ I=\frac{E}{R} \tag{②} \]

	电流沿顺时针方向。两金属杆都要受到安培力作用，作用于杆 $x_{1}y_{1}$ 的安培力为

	\[ F_{1}=BIl_{1} \tag{③} \]

	方向向上，作用于杆 $x_{2}y_{2}$ 的安培力为

	\[ F_{2}=BIl_{2} \tag{④} \]

	方向向下。当杆做匀速运动时，对系统受力分析得

	\[ F-m_{1}g-m_{2}g+F_{1}-F_{2}=0 \tag{⑤} \]

	联立解得

	\[ I=\frac{F-(m_{1}+m_{2})g}{B(l_{2}-l_{1})} \tag{⑥} \]

	\[ v=\frac{F-(m_{1}+m_{2})g}{B^{2}(l_{2}-l_{1})^{2}}R \tag{⑦} \]

	作用于两杆的重力的功率的大小

	\[ P=(m_{1}+m_{2})gv \tag{⑧} \]

	回路电阻上的热功率为

	\[ Q=I^{2}R \tag{⑨} \]

	由⑥⑦⑧⑨式，可得

	\[ P=\frac{F-(m_{1}+m_{2})g}{B^{2}(l_{2}-l_{1})^{2}}R(m_{1}+m_{2})g, \]

	\[ Q=\left[\frac{F-(m_{1}+m_{2})g}{B(l_{2}-l_{1})}\right]^{2}R. \]
}

\item
%% number: 25
%% paperName: 2004年全国理综Ⅰ第 25 题 20 分
%% typeId: 6
%% type: 计算题
%% chapter: 运动和力
%% point: 牛顿定律的应用
%% method: 无
%% score: 20
%% degree: 450
%% duplicateId: 0
%% body:
一小圆盘静止在桌布上，位于一方桌的水平面的中央。桌布的一边与桌的 AB 边重合，如图。已知盘与桌布间的动摩擦因数为 $\mu_{1}$，盘与桌面间的动摩擦因数为 $\mu_{2}$。现突然以恒定的加速度 $a$ 将桌布抽离桌面，加速度的方向是水平的且垂直于 AB 边。若圆盘最后未从桌面掉下，则加速度 $a$ 满足的条件是什么？（以 $g$ 表示重力加速度）
\onepicture[align=right, w=7.0cm]{\includesvg{figs/25}}
%% answer: 
\jdanswer{
\[ a\geq\frac{\mu_{1}+2\mu_{2}}{\mu_{2}}\mu_{1}g \]
}
%% memo:
\memoanswer{
	设圆盘的质量为 $m$，桌长为 $l$，在桌布从圆盘下抽出的过程中，圆盘的加速度为 $a_{1}$，由牛顿第二定律得

	\[ \mu_{1}mg=ma_{1}, \]

	桌布抽出后，圆盘在桌面上做匀减速运动，以 $a_{2}$ 表示加速度的大小，由牛顿第二定律得

	\[ \mu_{2}mg=ma_{2}, \]

	设圆盘刚离开桌布时的速度为 $v_{1}$，移动的距离为 $x_{1}$，离开桌布后在桌面上再运动距离 $x_{2}$ 后便停下，由运动学公式得

	\[ v_{1}^{2}=2a_{1}x_{1},\quad v_{1}^{2}=2a_{2}x_{2}, \]

	圆盘没有从桌面上掉下的条件是

	\[ x_{1}+x_{2}\leq\frac{1}{2}l, \]

	设桌布从盘下抽出所经历时间为 $t$，在这段时间内桌布移动的距离为 $x$，由运动学公式得

	\[ x=\frac{1}{2}at^{2},\quad x_{1}=\frac{1}{2}a_{1}t^{2}, \]

	又 $x=\frac{1}{2}l+x_{1}$，联立解得 $a=\frac{\mu_{1}+2\mu_{2}}{\mu_{2}}\mu_{1}g$，故圆盘最后未从桌面掉下，则加速度 $a$ 满足的条件为

	\[ a\geq\frac{\mu_{1}+2\mu_{2}}{\mu_{2}}\mu_{1}g. \]
}

\end{enumerate}

\end{document}
